\origpage{51}{N. Robertson, P.D. Seymour / Journal of Combinatorial Theory, Series B 89 (2003) 43--76}
\textit{Subproof.} Let $v \in X_i \cap X_k$. Then $v \in V(A_i \cap A_k)$, and hence $v \in V(A_j)$ by (d). But we may assume that $j \ne i$, and so $V(A_i) \subseteq V(B_j)$ by (1); and hence $v \in V(B_j)$, and consequently $v \in X_j$. This proves (3).

Since $|X_i| \le 2\rho + 1$ for $1 \le i \le n$, it follows from (3) that every torso of $(H, M)$ is a $(2\rho + 1)$-ring with perimeter $t_1, \dots, t_n$. This proves statement 5 of 3.2, and hence completes the proof of 3.2. $\square$

\textbf{Proof of 1.3} (assuming 3.1). Let $\theta, \kappa, \rho, \zeta$ be as in 3.1. We may assume (by replacing $\theta$ by a larger number if necessary) that $\theta \ge 2\rho + 4 + \zeta$. Let $r = 2\rho + 1$, $d = \kappa$ and $w = \zeta + 4\theta - 3$. We claim that $r, d, w$ satisfy 1.3. For let $\mathcal{D}$ be the class of all designs $(G', M)$ such that every torso of $(G', M)$ is isomorphic to an outgrowth by $\le d$ $r$-rings of a graph in one of $\Sigma_1, \dots, \Sigma_s$, where $\Sigma_1, \dots, \Sigma_s$ are the surfaces in which $L$ cannot be drawn (up to homeomorphism).

\hspace*{12pt}(1) \textit{For any graph $G'$ with no $L$-minor, $\mathcal{D}^\zeta$ is $\theta$-pervasive in $G'$.}

\textit{Subproof.} Let $G$ be a subgraph of $G'$, and let $\mathscr{T}$ be a tangle in $G$ of order $\ge \theta$. Certainly $\mathscr{T}$ controls no $L$-minor of $G$, and so, by 3.1, there exists $Z \subseteq V(G)$ with $|Z| \le \zeta$, and a segregation $\mathcal{S}$ of $G \setminus Z$ of type $(\rho, \kappa)$, such that

\hspace*{12pt}$\bullet$ there is a proper arrangement $\alpha$ of $\mathcal{S}$ in one of $\Sigma_1, \dots, \Sigma_s$, and $\bullet$ for each $(A, \Omega) \in \mathcal{S}$ there is no $(A', B') \in \mathscr{T}$ with $Z \subseteq V(A' \cap B')$ and $B' \setminus Z \subseteq A$.

Let $(A_1, \Omega_1), \dots, (A_k, \Omega_k)$ be the members $(A, \Omega)$ of $\mathcal{S}$ with $|\bar{\Omega}| > 3$. For $1 \le j \le k$, let the vertices of $\bar{\Omega}_j$ be $t_1^j, \dots, t_{n_j}^j$, in order. Since $(A_j, \Omega_j)$ is a $\rho$-vortex, by 3.2 there are separations $(A_1^j, B_1^j), \dots, (A_{n_j}^j, B_{n_j}^j)$ of $A_j$ as in 3.2. For each $(A, \Omega) \in \mathcal{S}$, let $C(A, \Omega)$ be the unique subgraph $B$ of $G \setminus Z$ such that $(A, B)$ is a separation of $G \setminus Z$ and $V(A \cap B) = \bar{\Omega}$. For $1 \le j \le k$ and $1 \le i \le n_j$, let $C_i^j = \delm{B}\fix{B_i^j} \cup C(A_j, \Omega_j)$.\fixnote{fix:1.3-Cij}{No $B$ is defined here, and the sets $B_i^j$ are those of the separations $(A_i^j, B_i^j)$ just chosen.} Let
\begin{equation*}
\mathcal{L} = \{(A_i^j, C_i^j) : 1 \le j \le k, 1 \le i \le n_j\} \cup \{(A, C(A, \Omega)) : (A, \Omega) \in \mathcal{S}, |\bar{\Omega}| \le 3\}.
\end{equation*}
Then since $\mathcal{S}$ is a segregation of $G$, and $(A_1^j, B_1^j), \dots, (A_{n_j}^j, B_{n_j}^j)$ satisfy 3.2, it is straightforward to verify \del{(we omit the details)} that $\mathcal{L}$ is a location in $G \setminus Z$ and every torso of the design of $\mathcal{L}$ is an outgrowth by $\le d$ $r$-rings of a graph in one of $\Sigma_1, \dots, \Sigma_s$, that is, $\mathcal{D}$ contains the design of $\mathcal{L}$. (We use here that for $(A, \Omega) \in \mathcal{S}$ with $|\bar{\Omega}| \le 3$, the edges for the torso of $\mathcal{L}$ with ends in $\bar{\Omega}$ may all be drawn within $\alpha(A, \Omega)$.) \fix{We give the details. For $(A, \Omega) \in \mathcal{S}$ with $|\bar{\Omega}| \le 3$ let $\mathcal{L}(A, \Omega) = \{(A, B_0)\}$, where $B_0$ is the graph with vertex set $\bar{\Omega}$ and no edges, and for $1 \le j \le k$ let $\mathcal{L}(A_j, \Omega_j) = \{(A_i^j, B_i^j) : 1 \le i \le n_j\}$. In each case $\mathcal{L}(A, \Omega)$ is a location in $A$ (by 3.2 for the vortices), and $\bar{\Omega} \subseteq V(B)$ for each $(A', B) \in \mathcal{L}(A, \Omega)$ (by statement 1 of 3.2). Thus $\mathcal{L}$ is the set of all $(A', B \cup C(A, \Omega))$ with $(A, \Omega) \in \mathcal{S}$ and $(A', B) \in \mathcal{L}(A, \Omega)$. Each of these is a separation of $G \setminus Z$, since $(A', B)$ is a separation of $A$ and $V(A \cap C(A, \Omega)) = \bar{\Omega} \subseteq V(B)$. Since $\mathcal{S}$ is a segregation, $C(A, \Omega)$ includes every other member of $\mathcal{S}$. Hence if two members of $\mathcal{L}$ come from distinct $(A, \Omega), (A^*, \Omega^*) \in \mathcal{S}$, the first side of the first is included in $A \subseteq C(A^*, \Omega^*)$; and if they come from the same $(A, \Omega)$, this holds since $\mathcal{L}(A, \Omega)$ is a location. So $\mathcal{L}$ is a location in $G \setminus Z$.

Now let $H$ be a torso of the design $(H_0, M)$ of $\mathcal{L}$. For $(A', B) \in \mathcal{L}(A, \Omega)$ we have $V(A' \cap (B \cup C(A, \Omega))) = V(A' \cap B)$, since $\bar{\Omega} \subseteq V(B)$. Assign each edge of $H$ not in $H_0$ to some $(A, \Omega) \in \mathcal{S}$ such that its ends lie in $V(A' \cap B)$ for some $(A', B) \in \mathcal{L}(A, \Omega)$, and each edge of $H_0$ to the member of $\mathcal{S}$ that contains it. For $(A, \Omega) \in \mathcal{S}$, let $P(A, \Omega)$ be the subgraph of $H$ with vertex set $V(H) \cap V(A)$ and with the edges assigned to $(A, \Omega)$. Since $A$ is included in every $C(A^*, \Omega^*)$ with $(A^*, \Omega^*) \ne (A, \Omega)$, a vertex of $A$ is in $H$ if and only if it is in $V(B)$ for every $(A', B) \in \mathcal{L}(A, \Omega)$, and the same holds for edges of $A$ and $H_0$. Hence $P(A, \Omega)$ is a torso of the design of $\mathcal{L}(A, \Omega)$. So $V(P(A, \Omega)) = \bar{\Omega}$ if $|\bar{\Omega}| \le 3$, and $P(A_j, \Omega_j)$ is a $(2\rho + 1)$-ring with perimeter $t_1^j, \dots, t_{n_j}^j$ by statement 5 of 3.2. Since $\mathcal{S}$ is a segregation, the graphs $P(A, \Omega)$ have union $H$, no two have a common edge, and $V(P(A, \Omega) \cap P(A^*, \Omega^*)) \subseteq \bar{\Omega} \cap \bar{\Omega}^*$. Also $V(\mathcal{S}) \subseteq V(H)$, since $\bar{\Omega} \subseteq V(B)$ for every $(A', B) \in \mathcal{L}(A, \Omega)$, and so $V(P(A, \Omega)) \cap V(\mathcal{S}) = \bar{\Omega}$.

Let $\alpha$ be the proper arrangement of $\mathcal{S}$ in $\Sigma \in \{\Sigma_1, \dots, \Sigma_s\}$, and let $G_0$ be the subgraph of $H$ with vertex set $V(\mathcal{S})$ and with the edges of all $P(A, \Omega)$ with $|\bar{\Omega}| \le 3$. Then $H = G_0 \cup P(A_1, \Omega_1) \cup \dots \cup P(A_k, \Omega_k)$. We draw $G_0$ in $\Sigma$ with each $x \in V(\mathcal{S})$ at $\alpha(x)$; for $|\bar{\Omega}| \le 3$ the edges of $P(A, \Omega)$ have ends in $\bar{\Omega}$, and since $|\bar{\Omega}| \le 3$ they may be drawn in $\alpha(A, \Omega)$ without crossings, meeting $bd(\alpha(A, \Omega))$ only in their ends. This is a drawing, since two of the discs meet only in points $\alpha(v)$ with $v \in V(\mathcal{S})$. Let $\Delta_j = \alpha(A_j, \Omega_j)$ $(1 \le j \le k)$. These discs are pairwise disjoint, since $\alpha$ is proper. If $x \in V(\mathcal{S})$ and $\alpha(x) \in \Delta_j$, then $x \in \bar{\Omega}_j$ by the second condition in the definition of an arrangement, and so $\alpha(x) \in bd(\Delta_j)$. Hence $\Delta_j$ meets the drawing of $G_0$ only in the points $\alpha(t_i^j)$, which lie on $bd(\Delta_j)$ in the order $t_1^j, \dots, t_{n_j}^j$. Finally $P(A_j, \Omega_j)$ meets $G_0$ just in $\bar{\Omega}_j$, since $V(P(A_j, \Omega_j)) \cap V(\mathcal{S}) = \bar{\Omega}_j$, and it is disjoint from $P(A_{j'}, \Omega_{j'})$ for $j' \ne j$, since $\Delta_j \cap \Delta_{j'} = \emptyset$. Thus $H$ is an outgrowth by $k \le \kappa = d$ $r$-rings of a graph in $\Sigma$.}\fixnote{fix:1.3-location-proof-added}{The paper omits the verification. The argument above is the one used in the formalization.}

For each $(A, B) \in \mathcal{L}$, let $(A^+, B^+)$ be the separation of $G$ such that $Z \subseteq V(A^+ \cap B^+)$, $A^+ \setminus Z = A$, $B^+ \setminus Z = B$ and $E(B^+)$ contains every edge of $E(G) \setminus E(A)$ with both ends in $V(B^+)$. Let $\mathcal{L}^+ = \{(A^+, B^+) : (A, B) \in \mathcal{L}\}$. Then $\mathcal{L}^+$ is a location in $G$, and its design belongs to $\mathcal{D}^\zeta$\fix{, once the separation $(K_e, G \setminus e)$, where $K_e$ consists of $e$ and its ends, is added to $\mathcal{L}^+$ for each edge $e$ of $G$ with an end in $Z$ and all its ends in $Z$ or in the core of $\mathcal{L}$ (such as an edge with both ends in $Z$, or a loop at a vertex of $Z$)}.\fixnote{fix:1.3-location}{A gap in the proof. Here the core of a location $\{(A_1, B_1), \dots, (A_k, B_k)\}$ in $G$ is the graph $G \cap B_1 \cap \dots \cap B_k$ of its design. An $n$-enlargement (page 47) requires every edge of the larger core to be an edge of the smaller one. Let $e$ be an edge of $G$ with an end in $Z$ and all its ends in $Z$ or in the core of $\mathcal{L}$. Then $e$ is in no $A$, since $A \subseteq G \setminus Z$, and its ends are in every $V(B^+)$, so $e$ is in every $B^+$. So $e$ is in the core of $\mathcal{L}^+$ but not in the core of $\mathcal{L}$, which lies in $G \setminus Z$. For example, if $G$ is one edge $uz$ with $Z = \{z\}$ and $\mathcal{L} = \emptyset$, the core of $\mathcal{L}^+$ contains the edge and the core of $\mathcal{L}$ does not. Adding $(K_e, G \setminus e)$ for each such $e$ keeps $\mathcal{L}^+$ a location, because $e$ is in every $B^+$ and in no $A^+$, and it removes these edges from the core. The added separations have order at most 2, so they are small by (2.7) of Graph Minors X. A new set $V(K_e)$ meets the core of $\mathcal{L}$ in at most one vertex $u$, and $u \in V(A \cap B)$ for some $(A, B) \in \mathcal{L}$, because the graphs $A$ with $(A, B) \in \mathcal{L}$ cover $G \setminus Z$ (for a vortex $(A_j, \Omega_j)$, by (a) in the proof of 3.2). So once the sets of the design of $\mathcal{L}$ are closed under subsets, with the empty set added, which does not change its torsos, the design of the new location is a $|Z|$-enlargement of it. So 1.3 holds as printed.} To complete the proof of (1) we must show that $\mathcal{L}^+ \subseteq \mathscr{T}$. Let $(A, B) \in \mathcal{L}$, and suppose that $(A^+, B^+) \notin \mathscr{T}$. Now $(A, B)$ has order $\le \max(3, 2\rho + 1)$ from the definition of $\mathcal{L}$, and so $(A^+, B^+)$ has order at most $\max(3, 2\rho + 1) + \zeta < \theta$. Consequently, $(B^+, A^+) \in \mathscr{T}$. Let $(A', B') = (B^+, A^+)$. Then
